The completion of this undergraduate thesis would not have been possible without the invaluable guidance, generous support, and enduring encouragement of numerous individuals, academic mentors, and public institutions.

First and foremost, the authors express their profound gratitude to the Almighty God, the source of all wisdom, perseverance, and strength, who guided our intellect and sustained our resolve throughout the demanding phases of this research endeavor.

We extend our sincere and deepest gratitude to our thesis adviser, whose intellectual rigor, constructive critiques, and steadfast guidance challenged us to pursue mathematical precision and academic excellence. Your insightful recommendations during the mathematical formulation of the Two-Stage MILP model, the linearization of non-linear deprivation functions, and the framing of distributive equity trade-offs were instrumental in shaping the scholarly quality of this manuscript.

We convey our heartfelt appreciation to the esteemed members of the Thesis Defense Panel and the faculty of the Department of Computer Science at the School of Engineering and Architecture, Ateneo de Davao University. Your incisive evaluations, meticulous methodological scrutiny, and substantive feedback during our milestone defenses immeasurably elevated the defensibility, operational relevance, and theoretical depth of this work.

We express our deepest and most sincere gratitude to the \textbf{Department of Science and Technology (DOST)} for funding and supporting this research undertaking. The financial grant, computational research resources, and invaluable institutional backing provided by DOST made this study possible and empowered our team to pursue rigorous computational research dedicated to enhancing disaster resilience and humanitarian logistics.

We gratefully acknowledge the governmental agencies and humanitarian organizations whose publicly accessible datasets, empirical frameworks, and technical assistance provided the empirical foundation for this study:
\begin{itemize}
    \item The \textbf{Department of Science and Technology (DOST)}, for generous research funding, technological support, and unwavering commitment to advancing scientific disaster risk reduction in the Philippines;
    \item The \textbf{Philippine Statistics Authority (PSA Region XI)}, for publishing the 2024 Census of Population (POPCEN) and the Philippine Standard Geographic Code (PSGC) records for Davao City's 182 administrative barangays;
    \item The \textbf{City Disaster Risk Reduction and Management Office (CDRRMO)} of Davao City, for providing operational context on municipal disaster response protocols, logistics staging hubs, and emergency dispatch procedures;
    \item The \textbf{Mines and Geosciences Bureau (MGB Region XI)}, for the flood and landslide geohazard susceptibility assessments that informed our spatial vulnerability modeling;
    \item The \textbf{Department of Social Welfare and Development (DSWD Field Office XI)}, for documenting regional relief stockpiling protocols and national disaster escalation guidelines;
    \item The \textbf{Sphere Association}, for establishing international humanitarian standards for minimum per-capita daily water and caloric requirements in acute crisis settings.
\end{itemize}

We also express our camaraderie and appreciation to our fellow Computer Science colleagues and research laboratory peers at Ateneo de Davao University, whose thought-provoking discussions, technical brainstorming, and shared experiences fostered an inspiring and collaborative research environment.

Finally, we dedicate this work with immense love and boundless gratitude to our families and parents. Your unwavering faith in our potential, countless sacrifices, moral encouragement, and unconditional love served as our pillar of strength and motivation throughout our academic journey. 

\vspace{2em}
\begin{flushright}
\textit{M. J. Misa} \\
\textit{K. L. Labor} \\
\textit{J. K. Gerona} \\
\vspace{1ex}
Ateneo de Davao University \\
October 2026
\end{flushright}
